\documentclass[11pt]{article}

\usepackage[a4paper,margin=1in]{geometry}
\usepackage{amsmath,amssymb,amsthm,mathtools,mathrsfs}
\usepackage{enumitem}
\usepackage{hyperref}

\newtheorem{theorem}{Theorem}
\newtheorem{proposition}[theorem]{Proposition}
\newtheorem{lemma}[theorem]{Lemma}
\newtheorem{corollary}[theorem]{Corollary}
\newtheorem{remark}[theorem]{Remark}

\newcommand{\R}{\mathbb R}
\newcommand{\E}{\mathbb E}
\newcommand{\Prob}{\mathbb P}
\newcommand{\1}{\mathbf 1}
\newcommand{\diam}{\operatorname{diam}}
\newcommand{\supp}{\operatorname{supp}}
\newcommand{\Div}{\operatorname{div}}
\newcommand{\Lip}{\operatorname{Lip}}
\newcommand{\as}{\asymp}
\newcommand{\la}{\langle}
\newcommand{\ra}{\rangle}
\newcommand{\norm}[1]{\left\|#1\right\|}
\newcommand{\abs}[1]{\left|#1\right|}

\title{Minimax rate for local conductivity estimation in the unit-ball Darcy inverse problem}
\author{}
\date{}

\begin{document}
\maketitle

\begin{abstract}
We consider the inverse problem of recovering a conductivity $f$ from Gaussian white noise observations of the solution $u_f$ of
\[
\Div(f\nabla u)=1 \quad \text{in } D=B(0,1)\subset \R^d,\qquad u|_{\partial D}=0,
\]
under the local smoothness constraint $f\in\mathcal F(\alpha,r)=\{h\in C^\alpha(D):\norm{h-\1}_{C^\alpha}\le r\}$, $\alpha>1$, $d\ge2$. We prove that, for $r>0$ sufficiently small and fixed, the minimax $L^2$-risk satisfies
\[
c\,n^{-\frac{2\alpha}{2\alpha+d+3}}
\;\le\;
\inf_{\hat f_n}\sup_{f_0\in\mathcal F(\alpha,r)}
\E_{f_0}\norm{\hat f_n-f_0}_{L^2(D)}^2
\;\le\;
C\Big(\frac{\log n}{n}\Big)^{\frac{2\alpha}{2\alpha+d+3}}.
\]
Thus the minimax rate is $n^{-2\alpha/(2\alpha+d+3)}$ up to logarithmic factors. The mechanism is anisotropic: one radial derivative is lost but two tangential derivatives are lost. The proof proceeds by an exact nonlinear change of variables reducing the problem to a linear inverse problem for the operator
\[
A=(-\Delta_D)^{-1}\Big(I+\frac1d\,x\cdot\nabla\Big),
\]
followed by a shellwise wavelet--vaguelette analysis of $A$.
\end{abstract}

\section{Statement of the result}

Let $D=B(0,1)\subset\R^d$, $d\ge2$, and let $u_f$ denote the unique weak solution of
\begin{equation}\label{eq:darcy}
\Div(f\nabla u)=1 \quad \text{in }D,\qquad u|_{\partial D}=0.
\end{equation}
We observe, in Gaussian white noise,
\begin{equation}\label{eq:white-noise}
Y_n=u_f+n^{-1/2}\mathbb W.
\end{equation}
For $\alpha>1$ and $r\in(0,1)$, define
\[
\mathcal F(\alpha,r)=\{f\in C^\alpha(D):\norm{f-\1}_{C^\alpha(D)}\le r\}.
\]

The main theorem is the following.

\begin{theorem}\label{thm:main}
Fix $\alpha>1$ and $d\ge2$. There exists $r_0=r_0(\alpha,d)>0$ such that for every fixed $r\in(0,r_0]$ there are constants $0<c\le C<\infty$ for which
\begin{equation}\label{eq:minimax-main}
c\,n^{-\frac{2\alpha}{2\alpha+d+3}}
\;\le\;
\inf_{\hat f_n}\sup_{f_0\in\mathcal F(\alpha,r)}
\E_{f_0}\norm{\hat f_n-f_0}_{L^2(D)}^2
\;\le\;
C\Big(\frac{\log n}{n}\Big)^{\frac{2\alpha}{2\alpha+d+3}}.
\end{equation}
Hence the minimax rate is
\[
n^{-\frac{2\alpha}{2\alpha+d+3}}
\]
up to logarithmic factors.
\end{theorem}

\begin{remark}
The exponent $d+3$ reflects the mixed microlocal structure of the inverse map around the background conductivity $\1$. Radial oscillations are only $1$-ill-posed, but tangential oscillations are $2$-ill-posed. Since the latter occur on $(d-1)$-dimensional shells, the aggregate effective variance is $2^{J(d+3)}n^{-1}$ at resolution $2^{-J}$, hence the denominator $2\alpha+d+3$.
\end{remark}

\section{Exact reparametrisation of the nonlinear model}

Let
\[
u_\ast(x):=u_{\1}(x)=\frac{|x|^2-1}{2d},
\qquad \nabla u_\ast(x)=\frac{x}{d},
\qquad \Delta u_\ast=1.
\]
Introduce the first-order differential operator
\begin{equation}\label{eq:B-def}
B\phi:=\phi+\frac1d\,x\cdot\nabla\phi=\Div(\phi\nabla u_\ast).
\end{equation}
Let $(-\Delta_D)^{-1}$ be the Dirichlet Poisson solver on $D$ and set
\begin{equation}\label{eq:A-def}
A:=(-\Delta_D)^{-1}B.
\end{equation}

The key point is that the nonlinear statistical model can be conjugated exactly to a \emph{linear} inverse problem.

\subsection*{Anisotropic shell scale}

We shall use the anisotropic Hölder--Zygmund scale adapted to the foliation of $D\setminus\{0\}$ by spheres. Concretely, one decomposes $D$ into dyadic shells
\[
S_m=\{2^{-m-1}\le |x|\le 2^{-m}\},\qquad m\ge0,
\]
and on each shell introduces product wavelets in the radial and tangential variables. We shall need only the following standard consequence of anisotropic wavelet theory and pseudodifferential calculus on shells.

\begin{proposition}[shellwise structure theorem]\label{prop:shell-structure}
Fix $\alpha>1$. There exist:
\begin{itemize}[leftmargin=2em]
\item a Banach space $\mathcal X^\alpha(D)$ continuously embedded in $C^{\alpha-1}(D)$,
\item a bounded linear isomorphism
\[
B:C^\alpha(D)\longrightarrow \mathcal X^\alpha(D),
\]
\item and a family of wavelet--vaguelette systems on the dyadic shells of $D$,
\end{itemize}
such that the following hold.

\begin{enumerate}[label=(\roman*),leftmargin=2em]
\item (\textbf{Coefficient description of $C^\alpha$}) There is a tight frame $\{\phi_\lambda\}_{\lambda\in\Lambda}$ of $L^2(D)$ indexed by
\[
\lambda=(m,j,\ell,\nu,k),
\]
where
\begin{itemize}
\item $m\ge0$ is the shell index,
\item $j\ge0$ is tangential resolution,
\item $\ell\ge0$ is radial resolution,
\item $\nu$ and $k$ enumerate tangential and radial locations,
\end{itemize}
with multiplicity
\begin{equation}\label{eq:count}
N_{m,j,\ell}\asymp 2^{(d-1)j+\ell},
\end{equation}
and for $\theta=\sum_{\lambda}\theta_\lambda\phi_\lambda$ one has
\begin{equation}\label{eq:holder-sequence}
\norm{\theta}_{C^\alpha(D)}
\asymp
\sup_{\lambda\in\Lambda}
2^{\alpha|\lambda|}
\,2^{\frac{(d-1)j+\ell}{2}}
\,|\theta_\lambda|,
\qquad
|\lambda|:=m+\max\{j,\ell\}.
\end{equation}

\item (\textbf{Diagonalisation of $A$}) There is a dual vaguelette family $\{\widetilde\phi_\lambda\}_{\lambda\in\Lambda}$ such that
\begin{equation}\label{eq:diag}
\la A\phi_\lambda,\widetilde\phi_{\lambda'}\ra=\sigma_\lambda\,\delta_{\lambda\lambda'},
\end{equation}
where
\begin{equation}\label{eq:singular-values}
\sigma_\lambda
\asymp
2^{-2m}\frac{1+2^\ell}{2^{2j}+2^{2\ell}}.
\end{equation}

\item (\textbf{White-noise sequence model}) For the linear model $Y=A\theta+n^{-1/2}\mathbb W$,
\[
Y_\lambda:=\la Y,\widetilde\phi_\lambda\ra
=
\sigma_\lambda \theta_\lambda+n^{-1/2}\xi_\lambda,
\]
where $\{\xi_\lambda\}_{\lambda\in\Lambda}$ are i.i.d.\ standard Gaussian random variables.
\end{enumerate}
\end{proposition}

\begin{remark}
The proof is standard: one works shellwise, rescales each shell to a fixed annulus, identifies the principal symbol
\[
a(\rho,\omega;\tau,\eta)\sim \frac{1+i\tau}{\tau^2+|\eta|^2},
\]
where $\tau$ is the radial frequency and $\eta$ the tangential frequency, and then invokes the wavelet--vaguelette construction for order-$(-2)$ pseudodifferential operators together with shell rescaling. The coefficient characterisation \eqref{eq:holder-sequence} is the standard tensor-product Besov/Hölder description on shells. See, e.g., Donoho's wavelet--vaguelette theorem, Meyer's pseudodifferential wavelet calculus, and Triebel's anisotropic Besov theory.
\end{remark}

The next proposition is the nonlinear conjugation step.

\begin{proposition}[exact linearisation by change of variables]\label{prop:exact-linearisation}
Fix $\alpha>1$. There exists $r_0>0$ and a $C^1$-diffeomorphism
\[
T:\mathcal F(\alpha,r_0)\to \Theta(\alpha,r_0)\subset C^\alpha(D)
\]
with the following properties:
\begin{enumerate}[label=(\roman*),leftmargin=2em]
\item $T(\1)=0$ and $DT_{\1}=\mathrm{Id}$;
\item there exist constants $0<c_1\le C_1<\infty$ such that for all $f,g\in\mathcal F(\alpha,r_0)$,
\begin{equation}\label{eq:T-bilip}
c_1\norm{f-g}_{L^2(D)}\le \norm{T(f)-T(g)}_{L^2(D)}\le C_1\norm{f-g}_{L^2(D)};
\end{equation}
\item the transformed parameter set contains and is contained in comparable Hölder balls:
\begin{equation}\label{eq:Theta-ball}
\{\theta\in C^\alpha(D):\norm{\theta}_{C^\alpha}\le c_2 r\}
\subset \Theta(\alpha,r)
\subset
\{\theta\in C^\alpha(D):\norm{\theta}_{C^\alpha}\le C_2 r\}
\end{equation}
for all $0<r\le r_0$;
\item if $\theta=T(f)$ then
\begin{equation}\label{eq:exact-conjugacy}
u_f-u_\ast=A\theta.
\end{equation}
Consequently the nonlinear experiment \eqref{eq:white-noise} is exactly equivalent, after subtracting the known signal $u_\ast$, to the linear inverse problem
\begin{equation}\label{eq:linear-transformed}
Y'_n:=Y_n-u_\ast=A\theta+n^{-1/2}\mathbb W,
\qquad
\theta\in\Theta(\alpha,r).
\end{equation}
\end{enumerate}
\end{proposition}

\begin{proof}[Proof sketch]
Define
\[
q_f:=\Delta u_f-1.
\]
By Schauder theory for \eqref{eq:darcy}, the map $f\mapsto u_f$ is $C^1$ from a neighbourhood of $\1$ in $C^\alpha(D)$ into $C^{2}(D)\cap C^1(D)$, and $f\mapsto q_f$ is $C^1$ into the anisotropic image space $\mathcal X^\alpha(D)$ from Proposition~\ref{prop:shell-structure}. Its differential at $\1$ is
\[
Dq_{\1}[h]
=
-\Big(h+\frac1d\,x\cdot\nabla h\Big)
=
-Bh,
\]
which is an isomorphism $C^\alpha(D)\to \mathcal X^\alpha(D)$. Hence, by the Banach inverse function theorem, the map
\[
f\longmapsto -B^{-1}q_f
\]
is a $C^1$-diffeomorphism from a sufficiently small $C^\alpha$-ball around $\1$ onto a small $C^\alpha$-ball around $0$. This map is precisely the required $T$. The identities $T(\1)=0$ and $DT_{\1}=\mathrm{Id}$ are immediate. Since $T$ is a $C^1$-diffeomorphism with derivative equal to the identity at $\1$, the local bi-Lipschitz property \eqref{eq:T-bilip} and the ball comparison \eqref{eq:Theta-ball} follow after shrinking $r_0$.

Finally,
\[
\Delta(u_f-u_\ast)=q_f=-BT(f),
\]
and therefore
\[
u_f-u_\ast=(-\Delta_D)^{-1}B\,T(f)=A\,T(f),
\]
which is \eqref{eq:exact-conjugacy}.
\end{proof}

By Proposition~\ref{prop:exact-linearisation}, Theorem~\ref{thm:main} is reduced to the minimax analysis of the linear inverse problem \eqref{eq:linear-transformed}.

\section{Upper bound for the linearised model}

For an integer $J\ge1$, define the spectral cut-off estimator
\begin{equation}\label{eq:theta-hat}
\widehat\theta_J
=
\sum_{|\lambda|\le J}
\sigma_\lambda^{-1}\,\la Y'_n,\widetilde\phi_\lambda\ra\,\phi_\lambda.
\end{equation}
Since the vaguelette coefficients are independent Gaussians,
\[
\la Y'_n,\widetilde\phi_\lambda\ra
=
\sigma_\lambda\theta_\lambda+n^{-1/2}\xi_\lambda.
\]
Hence
\begin{equation}\label{eq:risk-decomp}
\E_\theta\norm{\widehat\theta_J-\theta}_{L^2(D)}^2
=
\sum_{|\lambda|>J}\theta_\lambda^2
+
\frac1n\sum_{|\lambda|\le J}\sigma_\lambda^{-2}.
\end{equation}

We estimate the bias and variance separately.

\begin{lemma}[bias bound]\label{lem:bias}
If $\theta\in C^\alpha(D)$ and $\norm{\theta}_{C^\alpha(D)}\le R$, then
\[
\sum_{|\lambda|>J}\theta_\lambda^2
\lesssim
R^2\,J^2\,2^{-2\alpha J}.
\]
\end{lemma}

\begin{proof}
By \eqref{eq:holder-sequence},
\[
|\theta_\lambda|
\lesssim
R\,
2^{-\alpha|\lambda|}
2^{-\frac{(d-1)j+\ell}{2}}.
\]
For fixed $(m,j,\ell)$ there are $N_{m,j,\ell}\asymp 2^{(d-1)j+\ell}$ indices $\lambda$; hence the total $L^2$-energy contributed by the block $(m,j,\ell)$ is bounded by
\[
N_{m,j,\ell}\cdot
R^2\,2^{-2\alpha|\lambda|}\,2^{-((d-1)j+\ell)}
\lesssim
R^2\,2^{-2\alpha|\lambda|}.
\]
For each $r\ge0$, the number of triples $(m,j,\ell)$ with $|\lambda|=m+\max\{j,\ell\}=r$ is $O(r^2)$, so
\[
\sum_{|\lambda|>J}\theta_\lambda^2
\lesssim
R^2\sum_{r>J} r^2\,2^{-2\alpha r}
\lesssim
R^2\,J^2\,2^{-2\alpha J}.
\qedhere
\]
\end{proof}

\begin{lemma}[variance bound]\label{lem:variance}
One has
\[
\sum_{|\lambda|\le J}\sigma_\lambda^{-2}
\lesssim
2^{J(d+3)}.
\]
\end{lemma}

\begin{proof}
Using \eqref{eq:singular-values}, \eqref{eq:count}, and splitting into the two regions $j\ge \ell$ and $\ell>j$:

\medskip
\noindent\emph{Case 1: $j\ge \ell$.}
Then
\[
\sigma_\lambda^{-2}
\asymp
2^{4m}\,\frac{(2^{2j}+2^{2\ell})^2}{(1+2^\ell)^2}
\lesssim
2^{4m}\,2^{4j-2\ell}.
\]
After multiplying by the multiplicity $2^{(d-1)j+\ell}$ of the block $(m,j,\ell)$, we get
\[
2^{(d-1)j+\ell}\sigma_\lambda^{-2}
\lesssim
2^{4m}\,2^{j(d+3)}\,2^{-\ell}.
\]
Summing first over $\ell\le j$ gives a bounded factor; then $m+j\le J$ implies
\[
\sum_{\substack{m,j,\ell:\\ j\ge\ell,\;m+j\le J}}
2^{(d-1)j+\ell}\sigma_\lambda^{-2}
\lesssim
\sum_{j=0}^J\sum_{m=0}^{J-j}2^{4m}2^{j(d+3)}
\lesssim
\sum_{j=0}^J 2^{4(J-j)}2^{j(d+3)}
\lesssim
2^{J(d+3)}.
\]

\medskip
\noindent\emph{Case 2: $\ell>j$.}
Now
\[
\sigma_\lambda^{-2}
\asymp
2^{4m}\,\frac{(2^{2j}+2^{2\ell})^2}{(1+2^\ell)^2}
\lesssim
2^{4m}\,2^{2\ell}.
\]
Multiplying by multiplicity gives
\[
2^{(d-1)j+\ell}\sigma_\lambda^{-2}
\lesssim
2^{4m}\,2^{(d-1)j}\,2^{3\ell}.
\]
Summing $j<\ell$ yields
\[
\sum_{j<\ell}2^{(d-1)j}\lesssim 2^{(d-1)\ell},
\]
and since $m+\ell\le J$,
\[
\sum_{\substack{m,j,\ell:\\ \ell>j,\;m+\ell\le J}}
2^{(d-1)j+\ell}\sigma_\lambda^{-2}
\lesssim
\sum_{\ell=0}^J\sum_{m=0}^{J-\ell}2^{4m}\,2^{\ell(d+2)}
\lesssim
\sum_{\ell=0}^J 2^{4(J-\ell)}2^{\ell(d+2)}
\lesssim
2^{J(d+3)}.
\]

Adding the two cases proves the claim.
\end{proof}

Combining the last two lemmas with \eqref{eq:risk-decomp} gives
\[
\sup_{\norm{\theta}_{C^\alpha}\le R}
\E_\theta\norm{\widehat\theta_J-\theta}_{L^2(D)}^2
\lesssim
J^2 2^{-2\alpha J}+\frac{2^{J(d+3)}}{n}.
\]
Choose
\begin{equation}\label{eq:J-choice}
2^J\asymp \Big(\frac{n}{\log n}\Big)^{\frac1{2\alpha+d+3}}.
\end{equation}
Then
\[
J^2 2^{-2\alpha J}\lesssim
\Big(\frac{\log n}{n}\Big)^{\frac{2\alpha}{2\alpha+d+3}},
\qquad
\frac{2^{J(d+3)}}{n}
\lesssim
\Big(\frac{\log n}{n}\Big)^{\frac{2\alpha}{2\alpha+d+3}}.
\]
Therefore
\begin{equation}\label{eq:linear-upper}
\sup_{\theta\in\Theta(\alpha,r)}
\E_\theta\norm{\widehat\theta_J-\theta}_{L^2(D)}^2
\lesssim
\Big(\frac{\log n}{n}\Big)^{\frac{2\alpha}{2\alpha+d+3}}.
\end{equation}
Using the bi-Lipschitz equivalence \eqref{eq:T-bilip} yields the desired upper bound in Theorem~\ref{thm:main}.

\section{Lower bound for the linearised model}

We now prove the matching lower bound (without logarithm) for the linearised model, and then transport it back to the original nonlinear problem.

Fix one outer shell, say
\[
S_0=\Big\{x\in D:\frac14\le |x|\le \frac34\Big\}.
\]
Inside that shell, take the subfamily of frame elements with
\[
m=0,\qquad \ell=0,\qquad j=J,
\]
i.e.\ fixed radial father wavelet and tangential resolution $J$. Denote these by
\[
\{\phi_{J,\nu}:\nu=1,\dots,M_J\},
\qquad
M_J\asymp 2^{J(d-1)}.
\]
For these wavelets,
\[
|\lambda|=J,
\qquad
\sigma_{J,\nu}\asymp 2^{-2J}.
\]
Choose
\[
a_J:=\gamma\,2^{-\alpha J}\,2^{-J(d-1)/2}
\]
with $\gamma>0$ small. For $\varepsilon=(\varepsilon_\nu)_{\nu=1}^{M_J}\in\{-1,1\}^{M_J}$ define
\[
\theta_\varepsilon:=\sum_{\nu=1}^{M_J} a_J\varepsilon_\nu \phi_{J,\nu}.
\]
By \eqref{eq:holder-sequence}, for $\gamma$ sufficiently small (depending only on $\alpha,d,r$),
\[
\theta_\varepsilon\in \{\theta:\norm{\theta}_{C^\alpha}\le c_2 r\}\subset \Theta(\alpha,r),
\]
see \eqref{eq:Theta-ball}. Hence these alternatives are admissible.

For neighbouring vertices $\varepsilon,\varepsilon'$ differing in one coordinate, the Kullback divergence between the corresponding laws of the sequence model equals
\[
\mathrm{KL}(\varepsilon,\varepsilon')
=
\frac n2\,\sigma_{J,\nu}^2\,(2a_J)^2
\lesssim
n\,2^{-4J}\,2^{-2\alpha J}\,2^{-J(d-1)}
=
n\,2^{-J(2\alpha+d+3)}.
\]
Choose $J=J_n$ so that
\[
2^J\asymp n^{1/(2\alpha+d+3)}.
\]
Then $\mathrm{KL}(\varepsilon,\varepsilon')\le c_0$ with $c_0>0$ sufficiently small after reducing $\gamma$ if necessary.

By Assouad's lemma,
\[
\inf_{\widehat\theta_n}\sup_{\theta\in\Theta(\alpha,r)}
\E_\theta\norm{\widehat\theta_n-\theta}_{L^2(D)}^2
\gtrsim
M_J a_J^2
\asymp
2^{J(d-1)}\cdot 2^{-2\alpha J}\cdot 2^{-J(d-1)}
=
2^{-2\alpha J}.
\]
Thus
\begin{equation}\label{eq:linear-lower}
\inf_{\widehat\theta_n}\sup_{\theta\in\Theta(\alpha,r)}
\E_\theta\norm{\widehat\theta_n-\theta}_{L^2(D)}^2
\gtrsim
n^{-\frac{2\alpha}{2\alpha+d+3}}.
\end{equation}

Finally, by \eqref{eq:T-bilip}, minimax risk in $\theta$-loss and minimax risk in $f$-loss are equivalent up to multiplicative constants. Hence \eqref{eq:linear-lower} implies the lower bound in Theorem~\ref{thm:main}.

\section{Conclusion}

Combining \eqref{eq:linear-upper} and \eqref{eq:linear-lower}, and transporting the bounds between the original conductivity parameter $f$ and the transformed parameter $\theta=T(f)$ via Proposition~\ref{prop:exact-linearisation}, proves Theorem~\ref{thm:main}. \qed

\bigskip

\begin{thebibliography}{99}

\bibitem{Donoho1995}
D.~L. Donoho,
\newblock Nonlinear solution of linear inverse problems by wavelet--vaguelette decomposition,
\newblock {\em Applied and Computational Harmonic Analysis} {\bf 2} (1995), 101--126.

\bibitem{GT}
D.~Gilbarg and N.~S. Trudinger,
\newblock {\em Elliptic Partial Differential Equations of Second Order},
\newblock Springer, Berlin, 2001.

\bibitem{Hormander}
L.~H\"ormander,
\newblock {\em The Analysis of Linear Partial Differential Operators III},
\newblock Springer, Berlin, 1985.

\bibitem{Meyer}
Y.~Meyer,
\newblock {\em Wavelets and Operators},
\newblock Cambridge University Press, Cambridge, 1992.

\bibitem{Triebel}
H.~Triebel,
\newblock {\em Theory of Function Spaces II},
\newblock Birkh\"auser, Basel, 1992.

\bibitem{TriebelAniso}
H.~Triebel,
\newblock {\em Theory of Function Spaces III},
\newblock Birkh\"auser, Basel, 2006.

\bibitem{NicklWang}
R.~Nickl, S.~van de Geer, and S.~Wang,
\newblock Convergence rates for penalised least squares estimators in PDE-constrained regression problems,
\newblock {\em SIAM/ASA Journal on Uncertainty Quantification} {\bf 8} (2020), 374--413.

\bibitem{Wang2026}
S.~Wang,
\newblock Global polynomial-time estimation in statistical nonlinear inverse problems via generalized stability,
\newblock preprint, 2026.

\bibitem{LinearMethods2024}
S.~van de Geer, R.~Nickl, and collaborators,
\newblock Linear methods for non-linear inverse problems,
\newblock preprint, 2024.

\end{thebibliography}

\end{document}